% Page 026

\seite{26}

\abschnitt{Kaisergeburtstag}
\pstart
Die Kaisergeburtstagsfeier wurde gemeinschaftlich mit\ob
Schülern der Pfarrei in Seffern abgehalten. Es erließ\ob
bei einer alten Gepflogenheit. Die Festrede wurde von\ob
Herrn Lehrer Kiefrad aus Heilenbach gehalten.\ob
Montag daraufhin sind die drei Punkte: Habe\ob
und unserer Landesschriften

1.\ Lieben\ob
2.\ Ehren\ob
3.\ Sind Gehorsam sein.

Der Ortsschulinspector Herrn Pfarrer Ettges wurde
er gefragt gerecht, worauf die Verteilung der Kaiserw\ohb
ecken folgte.

\pend
\abschnitt{Osterprüfung}
\pstart

Die Osterprüfung wurde am 29.\ März durch den Orts\ohb
Schulinspektor abgehalten.

Am 4.\ April wurde die Schüler vom Herrn Kreisschul\ohb
inspektor revisiert. Dr. Reuter visitirt.

Mit dem Schlusse des Schuljahres und 6.\ April\ob
wurden fünf Schüler der Schule entlassen: Sämt\ohb
liche Schüler treten in das Geschäft bzw.\ in den\ob
Beruf des Vaters ein.

Mit Beginn des Schuljahres treten zwei dieser hiesiger\ob
Jungen in das Gymnasium zu Prüm ein, um sich\ob
zum Priesterberuf vorzubereiten. Privatim vorgebildet\ob
werden sie auf doch 4.\ Klasse (Untertertia) auf\ohb
genommen.
\pend
